\chapter{Potok Ricciego: metryka, krzywizna i pierwsze równania ewolucyjne}
\label{ch:potok-ricciego}
\index{potok Ricciego}

Metryka pozwala mierzyć długości i objętości, a tensor Ricciego
opisuje część jej krzywizny. Można więc zapytać, co się stanie,
gdy metrykę będziemy zmieniać z prędkością wyznaczoną przez
ten tensor. Na sferze oczekujemy kurczenia, ponieważ jej Ricci
jest dodatni; płaski torus powinien pozostać nieruchomy.
Oba przewidywania sprawdzimy rachunkiem. Następnie wyprowadzimy
równania dla miary i krzywizny skalarnej oraz wyjaśnimy, gdzie
w dowodzie istnienia rozwiązania potrzebna jest analiza równań
parabolicznych.

W całym rozdziale $M^n$ jest gładką, zwartą rozmaitością bez
brzegu, $n\geq2$, a metryki są riemannowskie. Orientacja nie jest
potrzebna do określenia gęstości objętości. Zachowujemy konwencję
krzywizny z Definicji~\ref{def:tensor-krzywizny}: okrągła sfera
ma dodatni tensor Ricciego. Piszemy
\[
 \Delta_g f=\operatorname{div}_g(\nabla^g f)
            =g^{ij}\nabla_i\nabla_j f.
\]
Jest to konwencja z rozdziału~\ref{ch:krzywizna}; na zwartej
rozmaitości $\int_M f\Delta_g f\,\dd\mu_g
=-\int_M|\nabla f|_g^2\,\dd\mu_g$. Jeżeli używamy
laplasjanu Hodge'a $\Delta_H=\dd\delta+\delta\dd$,
to na funkcjach $\Delta_H=-\Delta_g$. To rozróżnienie ustala
znak wszystkich poniższych równań cieplnych.

\section{Równanie i metryki o stałej krzywiźnie}
\label{sec:ricci-flow-definition}

\begin{definicja}[Potok Ricciego Hamiltona]
\label{def:ricci-flow}
Niech $g_0$ będzie gładką metryką na $M$. \emph{Potokiem
Ricciego} o wartości początkowej $g_0$ na przedziale $[0,T)$
nazywamy gładką rodzinę metryk $g(t)$ spełniającą
\begin{equation}\label{eq:ricci-flow}
 \partial_tg(t)=-2\operatorname{Ric}(g(t)),
 \qquad g(0)=g_0.
\end{equation}
Jest to wersja \emph{nienormalizowana}: nie dodajemy składnika,
który utrzymywałby całkowitą objętość na stałym poziomie.
\end{definicja}

Dla wektora $v\in T_pM$, trzymanego nieruchomo przy zmianie
czasu, równanie daje
\[
 \frac{\dd}{\dd t}|v|_{g(t)}^2
 =-2\operatorname{Ric}_{g(t)}(v,v).
\]
Dodatni Ricci zmniejsza więc kwadrat długości w danym
kierunku; ujemny go zwiększa. Nazwa „potok” oznacza tu
ewolucję tensora metrycznego, a nie przepływ punktów
rozmaitości z rozdziału~\ref{ch:pola}.

\begin{stwierdzenie}[Stała zmiana skali]
\label{prop:ricci-scaling}
Jeżeli $g_0$ jest metryką i $c>0$ jest liczbą niezależną
od punktu $M$, to dla $\widetilde g=cg_0$ zachodzi
\[
 \nabla^{\widetilde g}=\nabla^{g_0},\qquad
 \operatorname{Ric}(\widetilde g)=\operatorname{Ric}(g_0),
 \qquad \operatorname{Scal}(\widetilde g)
          =c^{-1}\operatorname{Scal}(g_0),
\]
\[
 K_{\widetilde g}=c^{-1}K_{g_0},\qquad
 \dd\mu_{\widetilde g}=c^{n/2}\dd\mu_{g_0}.
\]
Wzór na $K$ dotyczy każdej płaszczyzny stycznej.
\end{stwierdzenie}
\begin{proof}
W mapie $\widetilde g_{ij}=cg_{0,ij}$ i
$\widetilde g^{ij}=c^{-1}g_0^{ij}$. We wzorze na symbole
Christoffela~\eqref{eq:christoffel-metryka} czynnik $c$
z pochodnych metryki znosi $c^{-1}$ z macierzy
odwrotnej. Koneksja i operator $R(X,Y)Z$ są więc
niezmienione. Ricci z rozdziału~\ref{ch:krzywizna}
jest śladem endomorfizmu $W\mapsto R(W,X)Y$, zatem
również się nie zmienia jako tensor typu $(0,2)$.
Skalar wymaga jeszcze kontrakcji z
$\widetilde g^{ij}=c^{-1}g_0^{ij}$.

W ilorazie definiującym krzywiznę sekcyjną
\eqref{eq:krzywizna-sekcyjna} licznik jest liniowy
w metryce i mnoży się przez $c$, a wyznacznik
macierzy Grama w mianowniku przez $c^2$. Wreszcie
$\det(cg_{0,ij})=c^n\det(g_{0,ij})$; pierwiastek
wyznacznika daje ostatnią równość.
\end{proof}

\begin{stwierdzenie}[Metryka Einsteina jako rozwiązanie]
\label{prop:ricci-einstein-flow}
Załóżmy, że $\operatorname{Ric}(g_0)=\lambda g_0$
z pewną stałą $\lambda\in\R$. Wtedy
\begin{equation}\label{eq:ricci-einstein-solution}
 g(t)=(1-2\lambda t)g_0
\end{equation}
spełnia~\eqref{eq:ricci-flow} na każdym przedziale,
na którym $1-2\lambda t>0$. Gdy $\lambda>0$,
rozwiązanie tego kształtu kończy się przy
$t=(2\lambda)^{-1}$; gdy $\lambda=0$, jest stacjonarne;
gdy $\lambda<0$, skala rośnie.
\end{stwierdzenie}
\begin{proof}
Wstawmy $g(t)=c(t)g_0$ z $c(0)=1$. Ponieważ
$c(t)$ jest stałe na $M$, Stwierdzenie~\ref{prop:ricci-scaling}
daje $\operatorname{Ric}(g(t))=\lambda g_0$.
Równanie~\eqref{eq:ricci-flow} przyjmuje postać
$c'(t)g_0=-2\lambda g_0$, czyli $c'=-2\lambda$.
Całkując z warunkiem $c(0)=1$, dostajemy podany wzór.
Dodatniość $c(t)$ jest konieczna, by tensor był metryką.
\end{proof}

\begin{przyklad}[Okrągła sfera]
\label{prz:ricci-sphere}
Niech $g_{\mathrm{jedn}}$ będzie metryką sfery jednostkowej
$S^n$, a $g_0=a_0^2g_{\mathrm{jedn}}$, $a_0>0$.
Z~\eqref{eq:stala-krzywizna-tensor} i
Stwierdzenia~\ref{prop:ricci-scaling} otrzymujemy
\[
 \operatorname{Ric}(g_0)
 =(n-1)g_{\mathrm{jedn}}
 =\frac{n-1}{a_0^2}g_0.
\]
Zatem rozwiązanie z poprzedniego stwierdzenia ma postać
\begin{equation}\label{eq:ricci-sphere}
 g(t)=a(t)^2g_{\mathrm{jedn}},\qquad
 a(t)^2=a_0^2-2(n-1)t,
 \qquad 0\leq t<\frac{a_0^2}{2(n-1)}.
\end{equation}
Promień $a(t)$ maleje, a $K(t)=a(t)^{-2}$ rośnie.
W chwili granicznej tensor traci dodatnią określoność;
nie jest to gładkie przedłużenie potoku przez tę chwilę.
Dla $n=2$ pole powierzchni wynosi
$4\pi a(t)^2=4\pi a_0^2-8\pi t$.
\end{przyklad}

\begin{przyklad}[Płaski torus]
\label{prz:ricci-torus}
Niech $T^n=\R^n/\Lambda$ dla kraty $\Lambda$,
a $g_0$ niech będzie metryką pochodzącą od stałego
dodatnio określonego iloczynu skalarnego na $\R^n$.
W tych współrzędnych symbole Christoffela i tensor
krzywizny są zerowe, więc $\operatorname{Ric}(g_0)=0$.
Rodzina $g(t)=g_0$ spełnia~\eqref{eq:ricci-flow}.
W szczególności płaski torus dwuwymiarowy zachowuje
swój kształt i pole. Sam warunek $\chi(T^2)=0$ nie mówi,
że \emph{każda} metryka torusa jest płaska; dla innych
metryk pole pozostanie stałe, choć metryka może się
zmieniać.
\end{przyklad}

\section{Ewolucja miary i krzywizny skalarnej}
\label{sec:ricci-evolution}

\begin{stwierdzenie}[Zmiana elementu objętości]
\label{prop:ricci-volume}
Dla dowolnej gładkiej rodziny metryk o
$h=\partial_tg$ mamy
\[
 \partial_t\dd\mu_g
   =\tfrac12\operatorname{tr}_g h\,\dd\mu_g.
\]
Wzdłuż potoku Ricciego wynika stąd
\begin{equation}\label{eq:ricci-volume}
 \partial_t\dd\mu_g=-\operatorname{Scal}_g\,\dd\mu_g,
 \qquad
 \frac{\dd}{\dd t}\operatorname{Vol}_{g(t)}(M)
 =-\int_M\operatorname{Scal}_{g(t)}\,\dd\mu_{g(t)}.
\end{equation}
\end{stwierdzenie}
\begin{proof}
W mapie gęstość objętości to
$\sqrt{\det(g_{ij})}\,|\dd x^1\cdots\dd x^n|$.
Wzór Jacobiego na pochodną wyznacznika daje
\[
 \partial_t\det(g_{ij})
 =\det(g_{ij})\operatorname{tr}(g^{-1}h)
 =\det(g_{ij})g^{ij}h_{ij}.
\]
Pochodna pierwiastka wprowadza czynnik $1/2$.
Dla $h=-2\operatorname{Ric}$ ślad wynosi
$-2\operatorname{Scal}$. Zwarta rozmaitość pozwala
różniczkować całkę po $M$ pod znakiem całkowania.
\end{proof}

Do równania dla skalara potrzebujemy jednej tożsamości
wynikającej z krzywizny koneksji Leviego-Civity.

\begin{lemat}[Ściągnięta druga tożsamość Bianchiego]
\label{lem:ricci-contracted-bianchi}
Dla każdej gładkiej metryki riemannowskiej
\begin{equation}\label{eq:ricci-contracted-bianchi}
 \nabla^i\operatorname{Ric}_{ij}
       =\tfrac12\nabla_j\operatorname{Scal}.
\end{equation}
\end{lemat}
\begin{proof}
Użyjmy oznaczeń z~\eqref{eq:riemann-wspolrzedne}:
$R(\partial_i,\partial_j)\partial_b
=R^a{}_{b i j}\partial_a$ i
$\operatorname{Ric}_{bj}=R^a{}_{b a j}$; w drugim
wzorze użyliśmy symetrii Ricciego z
Rozdziału~\ref{ch:krzywizna}. Pokażmy krok prowadzący
do drugiej tożsamości Bianchiego. W punkcie wybieramy
mapę normalną i piszemy $D_i=\nabla_{\partial_i}$.
Ponieważ pola współrzędnych komutują,
$[D_i,D_j]V=R(\partial_i,\partial_j)V$.
Tożsamość Jacobiego dla operatorów $D_i$ ma postać
$[D_k,[D_i,D_j]]+[D_i,[D_j,D_k]]+[D_j,[D_k,D_i]]=0$.
Komutator $D_k$ z działaniem tensora $R_{ij}$
na $V$ jest $\nabla_kR_{ij}$, więc otrzymujemy
\emph{drugą tożsamość Bianchiego}
\[
 \nabla_kR^a{}_{b i j}
 +\nabla_iR^a{}_{b j k}
 +\nabla_jR^a{}_{b k i}=0.
\]
Równość jest tensorowa, więc obowiązuje w każdej
mapie. Kontrakcja $k=a$ daje
\[
 \nabla_aR^a{}_{b i j}
 -\nabla_i\operatorname{Ric}_{bj}
 +\nabla_j\operatorname{Ric}_{bi}=0.
\]
Po kontrakcji $b=i$ pierwszy składnik jest
$-\nabla^a\operatorname{Ric}_{aj}$, ponieważ symetrie
Riemanna dają $g^{bi}R^a{}_{b i j}
=-\operatorname{Ric}^a{}_j$.
Pozostałe składniki to odpowiednio
$-\nabla^i\operatorname{Ric}_{ij}$ i
$\nabla_j\operatorname{Scal}$. Ich suma jest zerem,
co daje~\eqref{eq:ricci-contracted-bianchi}.
\end{proof}

\begin{twierdzenie}[Ewolucja krzywizny skalarnej]
\label{thm:ricci-scalar-evolution}
Jeżeli $g(t)$ jest gładkim potokiem Ricciego, to
\begin{equation}\label{eq:ricci-scalar-evolution}
 \partial_t\operatorname{Scal}
 =\Delta_g\operatorname{Scal}
  +2|\operatorname{Ric}|_g^2
 =-\Delta_H\operatorname{Scal}
  +2|\operatorname{Ric}|_g^2.
\end{equation}
Drugą równość czytamy tylko w stopniu zero, gdzie
$\Delta_H=-\Delta_g$.
\end{twierdzenie}
\begin{proof}
Najpierw obliczmy zmianę skalara przy dowolnej
rodzinie $g(t)$ i oznaczmy $h_{ij}=\partial_tg_{ij}$.
Różniczkując $g^{ik}g_{kj}=\delta^i_j$, dostajemy
\begin{equation}\label{eq:ricci-inverse-variation}
 \partial_tg^{ij}=-g^{ip}h_{pq}g^{qj}=-h^{ij}.
\end{equation}
Różnica dwóch koneksji jest tensorem. Zróżniczkowanie
wzoru na koneksję Leviego-Civity i użycie $\nabla g=0$
dają
\begin{equation}\label{eq:ricci-connection-variation}
 \partial_t\Gamma^k_{ij}
 =\tfrac12g^{k\ell}
   (\nabla_i h_{j\ell}+\nabla_j h_{i\ell}
                        -\nabla_\ell h_{ij}).
\end{equation}
Z definicji Ricciego przez krzywiznę pochodna
składników kwadratowych w $\Gamma$ łączy się
z pochodnymi zwykłymi w pochodne kowariantne:
\[
 \partial_t\operatorname{Ric}_{ij}
 =\nabla_k(\partial_t\Gamma^k_{ij})
   -\nabla_i(\partial_t\Gamma^k_{kj}).
\]
Ponieważ $\partial_t\Gamma^k_{kj}
=\tfrac12\nabla_j(g^{pq}h_{pq})$, po kontrakcji
z $g^{ij}$ pierwszego składnika otrzymujemy
$\nabla_k\nabla_i h^{ik}
-\tfrac12\Delta_g(\operatorname{tr}_gh)$,
a z drugiego
$-\tfrac12\Delta_g(\operatorname{tr}_gh)$.
Różniczkując dodatkowo kontrakcję
$\operatorname{Scal}=g^{ij}\operatorname{Ric}_{ij}$
i korzystając z~\eqref{eq:ricci-inverse-variation},
otrzymujemy pełny wzór wariacyjny
\begin{equation}\label{eq:ricci-scalar-variation}
 \partial_t\operatorname{Scal}
 =-\langle h,\operatorname{Ric}\rangle_g
   +\nabla_k\nabla_i h^{ik}
   -\Delta_g(\operatorname{tr}_g h).
\end{equation}

Teraz $h=-2\operatorname{Ric}$, więc pierwszy składnik
jest $2|\operatorname{Ric}|^2$, a
$\operatorname{tr}_gh=-2\operatorname{Scal}$.
Lemat~\ref{lem:ricci-contracted-bianchi} daje
\[
 \nabla_k\nabla_i h^{ik}
 =-2\nabla_k\nabla_i\operatorname{Ric}^{ik}
 =-\Delta_g\operatorname{Scal}.
\]
Ostatni składnik~\eqref{eq:ricci-scalar-variation}
jest $+2\Delta_g\operatorname{Scal}$; suma obu
wyrazów laplasjanowych to $+\Delta_g\operatorname{Scal}$.
\end{proof}

\section{Powierzchnie, pole i zasada maksimum}
\label{sec:ricci-surfaces}

W wymiarze dwa mamy już
$\operatorname{Ric}=Kg$ i $\operatorname{Scal}=2K$
ze Stwierdzenia w Sekcji~\ref{sec:ricci-krzywizna-skalarna}.
Równanie potoku przyjmuje wtedy postać
\begin{equation}\label{eq:ricci-surface-metric}
 \partial_tg=-2Kg=-\operatorname{Scal}\,g.
\end{equation}
W danej chwili zmiana metryki jest więc proporcjonalna
do samej metryki: zachowuje lokalne kąty, choć współczynnik
skali zależy od punktu. Ponadto
$|\operatorname{Ric}|^2=|Kg|^2=2K^2$, ponieważ
$g^{ia}g^{jb}g_{ij}g_{ab}=2$. Dzieląc
\eqref{eq:ricci-scalar-evolution} przez dwa, dostajemy
\begin{equation}\label{eq:ricci-surface-curvature}
 \partial_tK=\Delta_gK+2K^2.
\end{equation}

\begin{stwierdzenie}[Pole zamkniętej powierzchni]
\label{prop:ricci-surface-area}
Niech $M^2$ będzie zwartą powierzchnią bez brzegu,
a $g(t)$ gładkim potokiem Ricciego. Wtedy
\begin{equation}\label{eq:ricci-surface-area}
 \operatorname{Area}_{g(t)}(M)
 =\operatorname{Area}_{g_0}(M)-4\pi\chi(M)t
\end{equation}
na całym przedziale istnienia rozwiązania.
\end{stwierdzenie}
\begin{proof}
Ze Stwierdzenia~\ref{prop:ricci-volume} i równości
$\operatorname{Scal}=2K$ wynika
\[
 \frac{\dd}{\dd t}\operatorname{Area}_{g(t)}(M)
 =-2\int_M K_{g(t)}\,\dd A_{g(t)}.
\]
Dopiero teraz stosujemy udowodzone dla powierzchni
zamkniętej twierdzenie Gaussa--Bonneta,
Wniosek~\ref{cor:gauss-bonnet-zamknieta}:
$\int_MK\,\dd A=2\pi\chi(M)$.
Zatem pochodna pola jest stałą $-4\pi\chi(M)$;
całkowanie po czasie daje~\eqref{eq:ricci-surface-area}.
Ten sam związek z charakterystyką Eulera wynika z
dwuwymiarowego przypadku twierdzenia Cherna--Gaussa--Bonneta
z rozdziału~\ref{ch:chern-weil}.
\end{proof}

Dla sfery $\chi(S^2)=2$, więc pole maleje z prędkością
$8\pi$, dokładnie jak w Przykładzie~\ref{prz:ricci-sphere}.
Dla torusa $\chi(T^2)=0$ pole jest stałe; metryka
płaska z Przykładu~\ref{prz:ricci-torus} jest szczególnym
rozwiązaniem stacjonarnym. Dla zwartej orientowalnej
powierzchni rodzaju $q>1$ mamy $\chi=2-2q<0$,
więc pole rośnie. Ostatni wniosek dotyczy pola każdego
istniejącego rozwiązania i nie stwierdza jego zbieżności.

\begin{twierdzenie}[Prosta zasada maksimum dla skalara]
\label{thm:ricci-scalar-max-principle}
Niech $g(t)$ będzie gładkim rozwiązaniem
\eqref{eq:ricci-flow} na zwartej rozmaitości bez brzegu,
na przedziale $[0,T)$. Wówczas funkcja
\[
 m(t)=\min_{x\in M}\operatorname{Scal}_{g(t)}(x)
\]
jest niemalejąca. W szczególności początkowa
nieujemność krzywizny skalarnej jest zachowana,
a jeśli $\operatorname{Scal}_{g_0}>0$ w każdym punkcie,
to pozostaje dodatnia w każdej chwili istnienia.
\end{twierdzenie}
\begin{proof}
Ustalmy $0\leq t_0<t_1<T$ i liczby $\varepsilon,\eta>0$.
Na $M\times[t_0,t_1]$ połóżmy
\[
 u(x,t)=\operatorname{Scal}_{g(t)}(x)
       -m(t_0)+\eta+\varepsilon(t-t_0).
\]
W chwili $t_0$ jest $u\geq\eta>0$.
Gdyby później $u$ osiągnęło zero, w pierwszej takiej
chwili $t_*>t_0$ istniałby punkt minimum $x_*$,
w którym $\partial_tu(x_*,t_*)\leq0$ i
$\Delta_{g(t_*)}u(x_*,t_*)\geq0$.
Jednak~\eqref{eq:ricci-scalar-evolution} daje tam
\[
 (\partial_t-\Delta_g)u
   =2|\operatorname{Ric}|^2+\varepsilon>0,
\]
sprzeczność. Zatem $u>0$ na całym cylindrze.
Przechodząc do granicy $\eta,\varepsilon\downarrow0$,
otrzymujemy
$\operatorname{Scal}_{g(t_1)}(x)\geq m(t_0)$
dla każdego $x$. Zwarta przestrzeń zapewnia
istnienie minimum; brak brzegu usuwa potrzebę
warunków brzegowych.
\end{proof}

\section{Poprawka DeTurcka i krótki czas istnienia}
\label{sec:ricci-deturck}

Równanie~\eqref{eq:ricci-flow} przypomina równanie
ciepła, bo Ricci zawiera drugie pochodne metryki.
Jego główny symbol ma jednak kierunki zerowe
pochodzące od swobody zmiany współrzędnych.
Pokażemy, jak dodatkowa pochodna Liego usuwa tę
degenerację. W tej sekcji istnienie rozwiązania
układu parabolicznego pozostaje jawnym wejściem
analitycznym.

Ustalmy pomocniczą gładką metrykę $\bar g$ na $M$,
na przykład $\bar g=g_0$, i oznaczmy jej symbole
Christoffela przez $\bar\Gamma^k_{pq}$. Różnica
koneksji jest tensorem, więc wzór
\begin{equation}\label{eq:ricci-deturck-vector}
 W^k(g)=g^{pq}\bigl(\Gamma^k_{pq}(g)
                            -\bar\Gamma^k_{pq}\bigr)
\end{equation}
określa globalne pole wektorowe. Składnik
$\bar\Gamma$ wskazuje punkt odniesienia dla
współrzędnych; nie zmienia krzywizny metryki $g$.
\emph{Równanie Ricciego--DeTurcka} ma postać
\begin{equation}\label{eq:ricci-deturck}
 \partial_t\widehat g
  =-2\operatorname{Ric}(\widehat g)
    +\Lie_{W(\widehat g)}\widehat g,
 \qquad \widehat g(0)=g_0.
\end{equation}

Wyprowadźmy część z najwyższymi pochodnymi
w lokalnej mapie. Z definicji symboli Christoffela
i krzywizny, pomijając tylko składniki zawierające
co najwyżej pierwsze pochodne $g$, otrzymujemy
\begin{equation}\label{eq:ricci-principal}
 (-2\operatorname{Ric}(g))_{ij}
 =g^{pq}\bigl(
   \partial_p\partial_q g_{ij}
  +\partial_i\partial_j g_{pq}
  -\partial_p\partial_i g_{qj}
  -\partial_p\partial_j g_{qi}\bigr)
  +\text{wyrazy niższego rzędu}.
\end{equation}
Z kolei po opuszczeniu indeksu pola~\eqref{eq:ricci-deturck-vector}
mamy
\[
 W_j=g^{pq}\bigl(\partial_p g_{qj}
                  -\tfrac12\partial_j g_{pq}\bigr)
      -g_{jk}g^{pq}\bar\Gamma^k_{pq}.
\]
Ponieważ $(\Lie_Wg)_{ij}=\nabla_iW_j+\nabla_jW_i$,
jego część drugiego rzędu to
\begin{equation}\label{eq:ricci-deturck-principal-correction}
 (\Lie_Wg)_{ij}
 =g^{pq}\bigl(
   \partial_i\partial_p g_{qj}
  +\partial_j\partial_p g_{qi}
  -\partial_i\partial_j g_{pq}\bigr)
  +\text{wyrazy niższego rzędu}.
\end{equation}
Pochodne ustalonej $\bar\Gamma$ i pochodne współczynników
$g^{pq}$ dają tylko wyrazy niższego rzędu względem
nieznanej metryki. Dodanie
\eqref{eq:ricci-principal} i
\eqref{eq:ricci-deturck-principal-correction} usuwa
wszystkie mieszane drugie pochodne. Równanie dla
każdej składowej ma więc część główną
\begin{equation}\label{eq:ricci-deturck-heat-principal}
 \partial_t\widehat g_{ij}
     =\widehat g^{pq}\partial_p\partial_q
             \widehat g_{ij}
       +\text{wyrazy niższego rzędu}.
\end{equation}
Po zamrożeniu współczynników w punkcie symbol
przestrzenny prawej strony to
$-g^{pq}\xi_p\xi_q\,\mathrm{Id}$ na tensorach
symetrycznych. Dodatnia określoność $g$ oznacza
$g^{pq}\xi_p\xi_q>0$ dla $\xi\ne0$: układ
\eqref{eq:ricci-deturck} jest ściśle paraboliczny.

\begin{twierdzenie}[Krótki czas istnienia; wejście analityczne]
\label{thm:ricci-short-time}
Dla każdej gładkiej metryki $g_0$ na gładkiej,
zwartej rozmaitości bez brzegu istnieje $T>0$
i dokładnie jedno gładkie rozwiązanie potoku Ricciego
z tą wartością początkową na dostatecznie krótkim
przedziale $[0,T)$. Jednoznaczność dotyczy gładkich
rozwiązań o tej samej wartości początkowej.
\end{twierdzenie}

\begin{proof}[Redukcja do twierdzenia parabolicznego]
\emph{Wejście analityczne} brzmi następująco: gładki
ściśle paraboliczny układ quasiliniowy drugiego rzędu
na zwartej rozmaitości bez brzegu, z gładką wartością
początkową, ma lokalnie w czasie jednoznaczne gładkie
rozwiązanie; ta sama lokalna teoria dotyczy równania
ciepła dla odwzorowań między zwartymi rozmaitościami,
również przy gładkiej metryce dziedziny zależnej od czasu.
Dowód tego twierdzenia wymaga oszacowań parabolicznych
i nie jest zawarty w niniejszym skrypcie. Dzięki
\eqref{eq:ricci-deturck-heat-principal} daje ono
rozwiązanie $\widehat g$ równania~\eqref{eq:ricci-deturck}.
Przez ciągłość $\widehat g(t)$ pozostaje dodatnio
określona dla dostatecznie małego $t$.

Niech $\phi_t$ rozwiązuje równanie zwyczajne
\[
 \partial_t\phi_t=-W(\widehat g(t))\circ\phi_t,
 \qquad \phi_0=\id_M.
\]
Gładkość pola i zwartość $M$ dają rodzinę
dyfeomorfizmów w krótkim czasie. Dla
$g(t)=\phi_t^*\widehat g(t)$ wzór na pochodną
cofnięcia tensora daje
\[
 \partial_tg
 =\phi_t^*\bigl(\partial_t\widehat g
                  -\Lie_{W(\widehat g)}\widehat g\bigr)
 =-2\phi_t^*\operatorname{Ric}(\widehat g)
 =-2\operatorname{Ric}(g).
\]
Ostatnia równość wynika z naturalności koneksji
Leviego-Civity i Ricciego względem dyfeomorfizmów.
To konstruuje rozwiązanie oryginalnego równania.

Dla jednoznaczności trzeba jeszcze przenieść
\emph{dowolne} rozwiązanie oryginalne do tego samego
układu parabolicznego. Dla gładkiego potoku $g(t)$
rozwiązujemy równanie ciepła dla odwzorowań
\[
 \partial_tF_t=\tau_{g(t),\bar g}(F_t),
 \qquad F_0=\id_M,
\]
gdzie $\tau$ jest śladem drugiej pochodnej
odwzorowania z $(M,g(t))$ do $(M,\bar g)$.
W mapach jego główna część to
$g^{ij}\partial_i\partial_jF^k$; lokalne istnienie
i jednoznaczność jest drugą częścią wskazanego
wejścia analitycznego. Przez krótki czas $F_t$ jest
dyfeomorfizmem, gdyż $F_0=\id_M$ i $M$ jest zwarte.
Metryka $\widehat g=(F_t^{-1})^*g$ spełnia
\eqref{eq:ricci-deturck}: po przeniesieniu przez
izometrię $F_t:(M,g)\to(M,\widehat g)$ pole
$\partial_tF_t\circ F_t^{-1}$ jest równe
$-W(\widehat g)$, gdyż napięcie identyczności
$(M,\widehat g)\to(M,\bar g)$ ma współrzędne
$\widehat g^{pq}(\bar\Gamma^k_{pq}
-\widehat\Gamma^k_{pq})$.
Dokładniej, dla $V_t=\partial_tF_t\circ F_t^{-1}$
wzór na pochodną cofnięcia daje
\[
 \partial_tg=F_t^*(\partial_t\widehat g
                       +\Lie_{V_t}\widehat g),
 \qquad V_t=-W(\widehat g).
\]
Stąd pochodna przeniesionej metryki zawiera
$+\Lie_{W(\widehat g)}\widehat g$.
Jednoznaczność układu parabolicznego wyznacza
$\widehat g$. Wówczas $F_t$ spełnia zwyczajne
równanie $\partial_tF_t=-W(\widehat g)\circ F_t$
z warunkiem $F_0=\id_M$, więc też jest wyznaczone
jednoznacznie; stąd $g=F_t^*\widehat g$.
Ten krok również
korzysta z podanego twierdzenia analitycznego;
rachunek symbolu sam w sobie nie dowodzi istnienia
ani jednoznaczności PDE.
\end{proof}

Analiza osobliwości, zbieżności i twierdzenia Perelmana
wymagają odrębnych narzędzi i nie wynikają z powyższego
twierdzenia o krótkim czasie.

\par\smallskip
{\footnotesize\noindent\emph{Dalsza lektura:}
\href{https://projecteuclid.org/journals/journal-of-differential-geometry/volume-17/issue-2/Three-manifolds-with-positive-Ricci-curvature/10.4310/jdg/1214436922.full}
{R.\ S.\ Hamilton, \emph{Three-manifolds with positive Ricci
curvature}, J. Differential Geom. 17 (1982), 255--306},
zwłaszcza \S\,3--4: pierwotne równanie i jego ewolucja;
\href{https://projecteuclid.org/journals/journal-of-differential-geometry/volume-18/issue-1/Deforming-metrics-in-the-direction-of-their-Ricci-tensors/10.4310/jdg/1214509286.full}
{D.\ M.\ DeTurck, \emph{Deforming metrics in the direction
of their Ricci tensors}, J. Differential Geom. 18 (1983),
157--162}: usunięcie degeneracji i redukcja do teorii
parabolicznej.\par}
